\chapter*{记号索引}
\phantomsection
\addcontentsline{toc}{chapter}{记号索引}

\(b^{J}, b^{J'} : 1, 2.\)

参考编号依次表示节与目（个别情形为习题）。

\(F^{J}\)（右模 F 的标量环 B 的反自同构）：1, 2。

\(N^{0}, M^{0}\)（N、M 为子模）：1, 3。

\(u^{*}\)（u 为同态）：1, 8。

\(M(x), \mathbf{x} (x \text{ 自由模中的元素}) : 1, 10.\)

\(M(u)\)（u 为自由模到自由模的同态）：1, 10。

\(^{t}M, \dot{M}^{J} (M \text { matrix } ): 1, 10.\)

\(\mathrm{D}_{\Phi}(x_{1},\ldots ,x_{n}),\mathrm{D}_{\Phi}(\mathrm{S}):2.\)

\(\overline{\alpha}\)（\(\alpha\) 为标量）：3。

\(\operatorname{Pf}(R):5,2.\)

\(\mathbf{Sp}(\Phi), \mathbf{Sp}(2m, \mathrm{A}), \mathbf{Sp}_{2m}(\mathrm{A}) : 5, 3.\)

\(\mathbf{U}(\Phi), \mathbf{S}\mathbf{U}(\Phi) (\Phi \text{ Hermitian 形式}) : 6, 2.\)

O(Q)、SO(Q)（Q 为二次型）：6, 2。

\(\mathbf{U}(n,\mathrm{A}),\mathbf{S}\mathbf{U}(n,\mathrm{A}),\mathbf{O}(n,\mathrm{A}),\mathbf{SO}(n,\mathrm{A}):6,2.\)

\(Q^{\top}Q^{\prime}, Q \sim Q^{\prime}(Q, Q^{\prime}\text{ 二次型}): 8, 1.\)

θ(Q)（Q 为二次型）：8, 1。

\(\mathbf{T} + \mathbf{T}', a. \mathbf{T} (\mathbf{T}, \mathbf{T}'\) 为二次型的型）：8, 2。

\(\delta (\mathbb{Q})\)（Q 为二次型）：8, 2。

\(\mathbf{Q} \otimes \mathbf{Q}'\)（Q、\(\mathbf{Q}'\) 为二次型）：8, 3。

\(\mathrm{TT}^{\prime}\)（T、\(\mathrm{T}^{\prime}\) 为二次型的型）：8, 3。

\(C(Q), I(Q), \rho_{Q}, \rho, C^{+}(Q), C^{-}(Q), C^{+}, C^{-} (Q \text{ 二次型}): 9, 1.\)

\(\alpha, \beta, \mathrm{G}(f): 9, 1.\)

\(e_x, i_f, i_x^{\mathrm{F}}, \lambda_{\mathrm{F}}: 9, 2.\)

\(\lambda_{\mathbf{F}}:9,3.\)

A(Φ)（Φ 为 2 维向量空间上的对称双线性型）：10, 1。

\(\overline{u} (u\) 为正相似变换）\(i,d:10,1\)。

\((\widehat{\mathbf{D}_1,\mathbf{D}_2})\) \((\mathbf{D}_1,\mathbf{D}_2\) 为直线或半直线）：10, 3。

\[
\mathfrak {A}, \mathfrak {A} _ {0}, h, h ^ {\prime}: 1 0, 3.
\]

\(|x|, \widehat{(x,y)}(x,y\text{ 向量}):10,3.\)

cos θ、sin θ、tan θ、cot θ：10, 3。

\(\left.\right)D_{1}, D_{2}\left\{\left(D_{1}, D_{2}\right)\left(D_{1}, D_{2} \text{ 定向平面中的半直线}\right): 10, 4.\right.\)

